% ═══════════════════════════════════════════════════════════════════════
%  Stratégie 2 --- XAUUSD niveaux x10. Chiffres : macros et tables générées.
% ═══════════════════════════════════════════════════════════════════════
\section{Anatomie d'un trade, pas à pas}
\label{app:x10_anatomie}

Cette annexe déroule la vie d'une transaction, de la barre où le niveau entre
dans le champ jusqu'à la barre où la position se ferme. Elle sert de clé de
lecture aux quatre exemples réels de la section~\ref{sec:x10_scenarios}.

\subsection{Les neuf étapes}

\begin{enumerate}
  \item \textbf{Le niveau est déterminé.} À chaque clôture de barre de cinq
  minutes, le moteur calcule les deux barreaux qui encadrent le prix. Seul
  celui vers lequel le prix se dirige est candidat.

  \item \textbf{La zone est testée.} Le prix doit être entré dans une bande
  de largeur $z \cdot ATR$ devant le barreau, sans l'avoir encore franchi.
  Trop loin, rien ne se passe~; déjà de l'autre côté, ce n'est plus une
  approche.

  \item \textbf{La cinématique est testée.} Vitesse vers le niveau au-dessus
  de son seuil, accélération vers le niveau au-dessus du sien, et déplacement
  de la dernière heure orienté vers le niveau. Si l'une des trois grandeurs
  est indéfinie --- amorçage d'indicateur, trou de données --- la barre est
  \emph{gelée}~: aucune transition n'a lieu, mais les compteurs continuent de
  courir.

  \item \textbf{L'automate arme.} Un seul niveau peut être armé à la fois~;
  tout autre niveau armable est ignoré. L'armement expire au bout de
  12~barres sans événement, ou dès que le prix ressort de la zone vers
  l'arrière.

  \item \textbf{Le niveau cède ou tient.} S'il cède --- clôture au-delà d'au
  moins un quart d'\code{ATR}, dans la moitié directionnelle de la barre ---
  on est sur un \emph{breakout}, et s'ouvre une période de maintien de trois
  barres. S'il tient --- débordement d'au moins trois dixièmes d'\code{ATR}
  puis retour en clôture dans les trois barres, avec décélération et
  basculement de l'accélération --- on est sur un \emph{sweep}.

  \item \textbf{La géométrie est posée.} Cible~: le barreau suivant dans le
  sens de la position. Stop~: de l'autre côté du niveau pour un
  \emph{breakout}, au-delà de l'extrême du débordement pour un
  \emph{reversal}. L'\code{ATR} utilisé est celui de la \emph{barre de
  décision}, jamais celui de la barre de cassure --- les deux diffèrent de
  trois barres.

  \item \textbf{Les quatre filtres passent, dans cet ordre.} Contexte, puis
  garde-fou gain/risque, puis fenêtre horaire --- évaluée sur la barre
  d'exécution, pas sur celle de décision --- puis taille. L'ordre n'est pas
  cosmétique~: il définit les dénominateurs des statistiques de refus
  publiées en section~\ref{sec:x10_grille}.

  \item \textbf{L'ordre est passé.} À l'ouverture de la barre suivante, et
  seulement si celle-ci commence exactement cinq minutes plus tard. Le prix
  payé est le prix médian augmenté ou diminué d'une demi-fourchette, selon le
  sens.

  \item \textbf{La position se résout.} Sur les barres d'une minute internes.
  Stop et cible sont testés minute par minute~; en cas de double contact dans
  la même minute, le stop l'emporte. À défaut, la position se ferme à la
  48\textsuperscript{e} barre de cinq minutes, ou à la clôture de séance.
\end{enumerate}

\subsection{Les quatre exemples, côte à côte}

% GÉNÉRÉ par scripts/build_x10_report_assets.py — ne pas éditer à la main.
% Toute valeur vient de results/xau_x10/*.json.
\begin{table}[htbp]
\centering
\scriptsize
\setlength{\tabcolsep}{4pt}
\caption{Les quatre trades détaillés en annexe. Règle de choix, écrite avant de regarder~: le trade de $R$ réalisé médian de chaque scénario parmi ceux clôturés en 2023.}
\label{tab:x10_trade_examples}
\begin{tabular}{l l r r r r r l r}
\toprule
\textbf{Scénario} & \textbf{Décision (UTC)} & \textbf{Niveau} & \textbf{Entrée} & \textbf{Stop} & \textbf{Cible} & \textbf{Sortie} & \textbf{Raison} & \textbf{R réalisé} \\
\midrule
Breakout Long & 2023-12-27 11:40 & $2\,080$ & $2\,081{,}800$ & $2\,078{,}209$ & $2\,090$ & $2\,078{,}209$ & STOP & $-1{,}00$ \\
Breakout Short & 2023-06-19 21:30 & $1\,950$ & $1\,948{,}735$ & $1\,950{,}822$ & $1\,940$ & $1\,950{,}822$ & STOP & $-1{,}00$ \\
Reversal Long & 2023-10-11 00:40 & $1\,860$ & $1\,860{,}175$ & $1\,859{,}630$ & $1\,870$ & $1\,859{,}630$ & STOP & $-1{,}00$ \\
Reversal Short & 2023-05-01 08:15 & $1\,990$ & $1\,988{,}940$ & $1\,992{,}261$ & $1\,980$ & $1\,992{,}261$ & STOP & $-1{,}00$ \\
\bottomrule
\end{tabular}
\end{table}


La règle de choix a été écrite dans le générateur avant d'être exécutée~: pour
chaque scénario, la transaction dont le $R$ réalisé est \emph{médian} parmi
celles clôturées en 2023, les égalités étant départagées par l'horodatage de
décision. Aucun tri esthétique n'intervient.

\begin{warningbox}
\textbf{Les quatre se terminent au stop, et c'est l'information.} Puisque
\XTenStopShare{} des transactions de la campagne sortent au stop --- pour une
espérance de \XTenStopR{}~R --- le \emph{trade médian} de n'importe laquelle
de ces populations est nécessairement un stop. Obtenir quatre exemples
gagnants aurait demandé une règle de sélection écrite après coup.

Cette remarque est aussi une leçon de lecture applicable ailleurs~: quand une
annexe d'exemples montre quatre transactions réussies, la question à poser est
« selon quelle règle ont-elles été choisies, et cette règle a-t-elle été
écrite avant~? ».
\end{warningbox}

\subsection{Lire une figure d'exemple}

Chaque figure de la section~\ref{sec:x10_scenarios} se lit de la même façon.

\begin{description}
  \item[Les barres.] Chandeliers de cinq minutes, du contexte qui précède la
  décision jusqu'à quelques barres après la sortie. Vert quand la clôture
  dépasse l'ouverture, rouge sinon.
  \item[La ligne foncée horizontale.] Le niveau x10 concerné.
  \item[La bande claire.] La zone d'approche, de largeur $z \cdot ATR$ devant
  le niveau. Elle est tracée à l'\code{ATR} de la transaction, donc sa hauteur
  varie d'un exemple à l'autre --- c'est la normalisation en action.
  \item[Les deux pointillés.] Le stop et la cible, tels qu'ils ont été posés à
  la barre de décision.
  \item[Le triangle montant.] L'entrée, au prix de remplissage réel, à
  l'ouverture de la barre qui suit la décision.
  \item[Le triangle descendant.] La sortie, au prix réel, étiquetée de sa
  raison.
\end{description}

\subsection{Deux cas de bord que les exemples ne montrent pas}

Ils sont rares mais normatifs, et un lecteur qui reproduirait la stratégie
doit les connaître.

\begin{description}
  \item[Le breakout raté qui devient reversal.] Si, pendant la période de
  maintien, une clôture repasse du mauvais côté du niveau, le \emph{breakout}
  est annulé et cette barre même devient candidate à un retournement --- sans
  nouvel armement, et en héritant de l'attribut de retest du \emph{breakout}
  avorté. Les deux familles de scénarios ne sont donc pas étanches.
  \item[Le sweep d'une seule barre.] La fenêtre de retour est inclusive et
  compte la barre du débordement elle-même. Un débordement en mèche suivi
  d'une clôture déjà revenue, dans la \emph{même} barre de cinq minutes,
  constitue un \emph{sweep} valide et donne lieu à une position. Conséquence
  assumée de la spécification.
\end{description}
